% year: תשע"ט
% semester: א
% moed: א
% number: 2
% points: 20
% categories: taylor
% source: exams/מבחנים/תשעט/מועד א תשעט.pdf

\begin{question}
(20 נק') כדי לחשב $a=\sqrt[3]{29}$ באופן מקורב מציגים
\[ a=3\cdot\sqrt[3]{29/27}=3\cdot\sqrt[3]{1+2/27} \]
ומשתמשים בפולינום מקלורין ממעלה 2 של פונקציה $f(x)=3\cdot(1+x)^{1/3}$. העריכו את שגיאת הקירוב בעזרת נוסחת טיילור.
\end{question}

\begin{hint}
השארית בצורת לגרנז' היא $R_2(x)=\frac{f'''(c)}{3!}x^3$ עבור $c$ בין $0$ ל-$x$. חסמו את $|f'''(c)|$ עבור $c\in(0,\frac2{27})$.
\end{hint}

\begin{solution}
\textbf{פולינום מקלורין.} נגזור את $f(x)=3(1+x)^{1/3}$:
\[ f'(x)=(1+x)^{-2/3},\qquad f''(x)=-\frac23(1+x)^{-5/3},\qquad f'''(x)=\frac{10}{9}(1+x)^{-8/3}. \]
לכן $f(0)=3$, $f'(0)=1$, $f''(0)=-\frac23$ ו-
\[ P_2(x)=3+x-\frac13x^2. \]
\textbf{הקירוב.} עבור $x=\frac2{27}$:
\[ a\approx P_2\!\left(\tfrac2{27}\right)=3+\frac2{27}-\frac13\cdot\frac{4}{729}=3+\frac{162}{2187}-\frac{4}{2187}=\frac{6719}{2187}\approx3.072245. \]
\textbf{הערכת השגיאה.} לפי נוסחת טיילור עם שארית לגרנז', קיים $c\in(0,\frac2{27})$ כך ש-
\[ a-P_2\!\left(\tfrac2{27}\right)=R_2\!\left(\tfrac2{27}\right)=\frac{f'''(c)}{3!}\left(\frac2{27}\right)^3=\frac{10}{54}(1+c)^{-8/3}\cdot\frac{8}{19683}. \]
מכיוון ש-$c>0$ מתקיים $0<(1+c)^{-8/3}<1$, ולכן
\[ 0<R_2<\frac{10}{54}\cdot\frac{8}{19683}=\frac{40}{531441}\approx7.53\cdot10^{-5}. \]
כלומר השגיאה קטנה מ-$7.6\cdot10^{-5}$ (ובפרט מ-$10^{-4}$), והקירוב הוא קירוב מלמטה (השארית חיובית). לשם השוואה, $\sqrt[3]{29}=3.0723168\ldots$, והשגיאה בפועל היא כ-$7.17\cdot10^{-5}$.
\end{solution}
